\chapter{系统总体设计}

\section{设计目标与原则}

根据前述需求分析，本文将平台总体设计目标归纳为以下几个方面：一是保证自动评测结果的正确性、稳定性和可重复性；二是通过容器化与服务化设计提升平台在高并发场景下的资源利用率与扩展能力；三是在不影响客观判题的前提下引入 AI 辅助能力，提高反馈质量和学习支持效果；四是为后续模块扩展、系统维护和持续优化预留良好的工程接口。

围绕上述目标，系统设计遵循以下原则：其一，判题核心优先原则，即任何智能辅助功能都不得替代标准测试用例判题流程；其二，服务解耦原则，将用户交互、业务管理、任务调度、判题执行和 AI 反馈等能力进行相对独立划分；其三，安全可控原则，对代码执行环境实施严格隔离与资源约束；其四，可扩展原则，为后续接入更多语言、更多题型或更多教学功能保留接口空间。

\section{总体架构设计}

综合平台业务特点与技术选型方向，本文采用前后端分离、服务协同和容器化执行相结合的总体架构。系统可以划分为表现层、业务服务层、评测执行层、AI 辅助层和数据支撑层五个部分。表现层负责为学生、教师和管理员提供统一访问入口；业务服务层负责处理用户管理、题库管理、提交管理和结果展示等核心业务；评测执行层负责从任务队列中拉取评测任务，并在隔离环境中完成编译、运行和结果回传；AI 辅助层负责在判题完成后生成解释性反馈；数据支撑层则承担数据持久化、日志存储和系统状态记录等职责。

在该架构中，提交请求首先进入业务服务层，由提交管理模块生成评测任务并写入任务队列；随后评测执行层从队列中消费任务，在容器环境中执行用户程序，并将运行结果、日志和状态信息回传；业务服务层再根据结果更新提交记录并返回前端展示。若启用 AI 辅助功能，则业务服务层可在客观判题完成后将必要信息发送给 AI 服务，以获得适度的错误解释和学习建议。整个流程体现了“判题主线确定、辅助模块后置”的设计思路。

\section{功能模块设计}

从功能划分角度看，平台主要包括用户管理模块、题库管理模块、提交与评测模块、结果与统计模块、AI 辅助反馈模块以及系统管理模块。用户管理模块负责身份认证、角色权限和账户信息维护；题库管理模块负责题目内容维护、样例输入输出配置和测试数据组织；提交与评测模块负责代码接收、任务入队、判题执行和结果生成；结果与统计模块负责提交记录展示、通过率统计和过程数据汇总；AI 辅助反馈模块负责生成错误分析、代码解释和优化建议；系统管理模块则负责服务配置、监控日志与资源维护。

其中，提交与评测模块是整个平台的核心。该模块既要完成业务层面的提交状态管理，又要与底层评测执行环境协同工作，因此应采用异步任务驱动模式。AI 辅助反馈模块则是平台的增强能力，它在设计上不进入客观判分链路，而是在主判题流程结束后以附加服务形式输出结果，从而保证系统架构在逻辑上清晰、在责任边界上明确。

\section{核心业务流程设计}

平台的核心业务流程可以概括为“提交受理、任务调度、隔离执行、结果回传、辅助反馈、统一展示”六个阶段。首先，学生通过前端提交代码，系统完成参数校验并生成提交记录；其次，提交记录转化为可调度任务进入队列；随后，评测执行节点获取任务并启动对应语言环境与资源约束，完成编译和用例执行；之后，判题结果被写回业务服务层并更新状态；若判题结果存在错误或用户主动请求解释，则调用 AI 辅助服务生成补充说明；最后，前端统一展示客观结果与可选的解释性反馈。

为了保证业务流程的稳定性，系统在设计上还需要考虑异常处理机制。例如，当编译阶段失败时，任务可直接返回编译错误结果而不进入运行阶段；当容器启动异常或执行超时时，系统应记录异常日志并返回相应状态；当 AI 服务不可用时，平台应仍能正常完成客观判题，只是暂时不提供智能辅助反馈。这样的异常兜底机制有助于避免辅助模块影响主流程稳定性。

\section{部署与运行思路}

在部署层面，平台适合采用容器化封装与分模块部署的方式。业务服务、评测工作节点和 AI 服务可以分别封装为独立容器镜像，在统一编排环境下运行。对于评测节点，可根据提交压力进行水平扩展；对于业务服务，可通过负载均衡提升访问稳定性；对于数据库与日志系统，则应关注持久化存储和备份恢复。这样的部署方式能够兼顾教学平台日常使用和高峰期任务处理需求。

在运行治理方面，平台应具备基本的可观测能力，包括任务执行日志、接口访问日志、服务运行状态、资源使用情况和异常告警信息等。借助统一日志记录与监控指标，系统可以更方便地完成性能分析、故障定位和容量评估。对在线评测平台而言，可观测性不仅关系到系统运维效率，也关系到教学场景下问题追踪和结果复盘的便利性。

\section{本章小结}

本章在需求分析基础上给出了平台的总体设计思路，明确了系统设计目标与原则，提出了面向在线评测场景的总体架构、功能模块划分、核心业务流程以及部署运行思路。该设计方案为后续系统详细实现、测试分析与论文完善提供了整体框架。
